\section{社会秩序犯罪 - 毒品犯罪}
\begin{enumerate}
  \item 走私、贩卖、运输、制造毒品。
        \begin{enumerate}
          \item 走私毒品。既遂标准，看陆（过了边境线）海空（到岸）。
          \item 贩卖毒品。没有购买毒品罪。
                \begin{enumerate}
                  \item 实行行为：卖出。
                  \item 既遂：\zd{卖掉毒品，不要求收到货款}。
                  \item 代购毒品。
                        \begin{enumerate}
                          \item 为贩毒者代购 \yjt 帮助犯；
                          \item 为吸毒者 \yjt\, \yi 有没有\zd{牟利目的}（中间商赚差价）；\er 不收钱分货（有没有\zd{贩卖目的}，自用不算）。
                        \end{enumerate}
                \end{enumerate}
          \item 运输毒品。
                \begin{enumerate}
                  \item 原来：运输是为了卖掉；\zd{现在不管目的}。\al{为了自己吸，出差携带毒品}。
                  \item 着手：为运输搬运；既遂：离开原存地，\zd{较长距离的位移（1km）}。
                  \item 和非法持有毒品的关系。\al{买来自己吸，没吸完，放家里，也算持有} \yjt 无缝打击。
                \end{enumerate}
          \item 制造毒品。\zd{既遂：把真正的毒品制造出来，不要求卖出去}；手段不能犯。
          \item 数量问题。
          \item 特别再犯制度。
        \end{enumerate}
  \item 非法持有毒品。轻罪。
  \item 容留他人吸毒。
        \begin{enumerate}
          \item 容留。给他人提供场所；
          \item 不能看名义，有人出钱有人出身份证。\al{甲乙开房吸毒，甲名义开房}。共有场所（大学宿舍），不构成容留。
        \end{enumerate}
  \item 引诱、教唆、欺骗他人吸毒和强迫他人吸毒。
  \item 包庇毒品犯罪分子。
  \item 窝藏、转移、隐瞒毒品、毒赃。
\end{enumerate}
